\documentclass[11pt]{article}
\usepackage[margin=1in]{geometry}
\usepackage{amsmath,amssymb,amsthm}
\usepackage{booktabs}
\usepackage{algorithm}
\usepackage[noend]{algpseudocode}
\usepackage[hidelinks]{hyperref}
\newtheorem{lemma}{Lemma}
\newtheorem{theorem}{Theorem}
\theoremstyle{remark}
\newtheorem{remark}{Remark}

\title{A Stable Samplesort with a Bounded Buffer:\\
Reusing Read Pages and Recording Flush Order}
\author{Lucas Sotomayor\\\small\texttt{github.com/sotomayorlucas/sotsort}}
\date{November 2026 \quad (draft)}
\renewcommand{\lg}{\log_2}

\begin{document}
\maketitle

\begin{abstract}
We describe a $k$-way distribution pass that is \emph{stable}, uses a buffer of
$B$ elements plus $O(n/P)$ words of tables (with pages of $P\approx B/k$
elements), and leaves every group contiguous in the input array after the
pass. Output pages live in the buffer; a full page is flushed to the next page
of the input that has already been read, and the page's group is recorded.
Because pages are always flushed to the next slot, slot index equals flush
order, which is what makes the pass stable. A permutation of slots by cycles
and a shift of the partial pages finish the pass. With sampled splitters this
gives a stable MSD samplesort. A pass costs $m\lceil\lg k\rceil$ comparisons and
at most $4.5m+O(kP)$ element moves. Inside SotSort, with the same memory
budget that VM-Powersort reserves, it sorts random \texttt{int}s with
$18.96$ comparisons and $22.10$ moves per element at $n\approx10^6$, against
$27.76$ and $31.65$ for VM-Powersort. The reuse of consumed blocks with recorded
order is not new (Polychroniou and Ross; Radsort, for LSD radix sort); what we
did not find in the literature is its use for a comparison-based, MSD,
stable distribution with contiguous groups. All counts are machine-independent
and reproducible from the repository.
\end{abstract}

\section{Problem}

A stable sort with a buffer of $B\ll n$ elements usually does $\Theta(\lg(n/B))$
binary passes over the data (stable quicksort with block partitions, or
mergesort with a bounded buffer). In SotSort's binary stable partition each
level costs about $1.5+c\lg(m/B)$ moves per element, with $c\approx0.56$--$0.75$
measured, because blocks partitioned in the buffer must be gathered with
rotations. Reducing the \emph{number of levels} is the lever. A $k$-way
distribution does the work of $\lg k$ binary levels in one pass. The unstable
version is standard (IPS$^4$o~\cite{ips4o}); the question is how to keep it
stable without $\Theta(n)$ extra space.

\section{The distribution pass}

Let the range be $a[0,m)$, with $k$ groups and pages of $P$ elements.
The buffer holds $k$ output pages and the $k-1$ splitters. A \emph{slot} is
$a[sP,(s+1)P)$. Let $w$ count flushed pages and $f_g$ the fill of the output
page of group $g$. The group of $x$ is the number of splitters $\le x$, so equal
keys fall in the same group.

\begin{algorithm}[h]
\caption{Stable paged distribution of $a[0,m)$ into $k$ groups}
\begin{algorithmic}[1]
\State $w\gets0$;\quad $f_g\gets0$ for all $g$
\For{$r=0,\dots,m-1$}
  \State $g\gets\textsc{Group}(a[r])$;\quad move $a[r]$ to page $g$ of the buffer;\quad $f_g\gets f_g+1$
  \If{$f_g=P$}
     \State move page $g$ to slot $w$;\quad $\mathit{grp}[w]\gets g$;\quad $w\gets w+1$;\quad $f_g\gets0$
  \EndIf
\EndFor
\State \textbf{Phase 1:} move the partial pages, in group order, to $a[wP,m)$
\State \textbf{Phase 2:} permute slots $0..w-1$ by group, stably, by cycles (one page of buffer)
\State \textbf{Phase 3:} move the partials to the buffer; for $g=k-1$ down to $0$, shift the
       slots of group $g$ right by $D_g=\sum_{h<g}f_h$ and put the partial of $g$ after them
\end{algorithmic}
\end{algorithm}

\begin{lemma}[a free slot exists]\label{l:free}
When a page is flushed with $w=s$, slot $s$ has been read completely and holds no flushed page.
\end{lemma}
\begin{proof}
At that moment $q=sP+\sum_g f_g$ elements have been read, and some $f_g=P$, so
$q\ge(s+1)P$. Slots $0..s-1$ hold the flushed pages; slot $s$ lies in the read
prefix and is not one of them.
\end{proof}

\begin{lemma}[slot index is flush order]\label{l:order}
Within each group, the stable order of its elements is: its flushed pages by
increasing slot index, then its partial page.
\end{lemma}
\begin{proof}
Each output page receives its elements in reading order. Pages are flushed to
slot $w$ and $w$ grows by one, so a page flushed earlier has a smaller index.
The partial page holds the last elements read for that group.
\end{proof}

\begin{lemma}[final free space]\label{l:space}
At the end, $a[wP,m)$ has exactly $\sum_g f_g$ positions, all read and none flushed.
\end{lemma}
\begin{proof}
$m=wP+\sum_g f_g$.
\end{proof}

\begin{theorem}
After phases 1--3, group $g$ occupies $[F_gP+D_g,\;F_{g+1}P+D_{g+1})$, where
$F_g$ is the number of full pages of groups $<g$, and the order within each
group is the input order.
\end{theorem}
\begin{proof}
Phase 1 is possible by Lemma~\ref{l:space}. In phase 2 a slot of group $g$ goes
to slot $F_g+e$, where $e$ is the number of earlier slots of group $g$; this is a bijection on
$0..w-1$ that preserves the order of slots inside a group, and it is applied by
following cycles with one page of buffer. In phase 3, since $D_g\le D_{g+1}$,
moving groups from right to left never overwrites anything not yet moved. The
order claim is Lemma~\ref{l:order}.
\end{proof}

Only the group of each slot must be recorded ($\lceil\lg k\rceil$ bits); our
implementation uses 2 bytes for the group and 4 for the cycle source, i.e.\ 6 bytes
per slot, $O(n/P)$ in total.

\paragraph{Cost of a pass.}
Comparisons: $m\lceil\lg k\rceil$ (for $k$ a power of two, exactly $m\lg k$).
Moves: $m$ to the buffer, $wP$ to flush, $\sum f$ in phase 1, at most $(M+C)P$ in
phase 2 with $M\le w$ misplaced slots and $C\le M/2$ cycles, then $\sum f$ twice and
at most $m$ for the shift in phase 3. Total $\le4m+3kP$ in the typical case and
$\le4.5m+3kP$ in the worst case of phase 2. One pass replaces $\lg k$ binary
levels; with $k=64$ that is about $0.7$ moves per binary level.

\paragraph{Classification without branches.}
The $k-1$ splitters form an implicit tree (node $i$ has children $2i$ and $2i+1$;
an in-order walk gives sorted order), and the group is $b-k$ after $\lg k$ steps of
$b\gets2b+[\text{splitter}_b\le x]$, as in IPS$^4$o's decision tree~\cite{ips4o}.
This changes no count.

\section{The sort}

Splitters: $8k$ random samples, sorted in the page area; $k-1$ equidistant ones
go to the tree. Ranges of size $\le B$ go to the base case, SotSort's stable
quicksort. Two safeguards bound bad inputs:
\begin{itemize}
\item a pass is \emph{bad} if its largest group exceeds $m/2$; two bad passes in a row on a
recursion path send the range to the base case; a pass that makes no progress
(one group holds the whole range) does so immediately;
\item depth is capped at $2\log_k n+2$.
\end{itemize}
The base case tolerates at most $4$ bad partitions per path before its merge
fallback (it previously tolerated $\lg n$). Hence the worst case stays
$O(n\log n)$ comparisons and moves; the constant is measured below, not proved tight.

In SotSort (version 8-r6) the pass is used only with a bounded block buffer
($2B<n$), for $k\in\{64,32,16,8\}$, the largest with $P\ge64$ and tables within
the plan's table reserve. Every byte is charged to the same budget as the rest
of the sort.

\section{Counts}

A \emph{move} is a move-assignment or move-construction of an element;
comparisons are calls to the comparator. Counts come from the benchmark
harness of VM-Powersort~\cite{vmps}, with SotSort given the memory VM-Powersort reserves for the
same input; the harness verifies every output. $S$ is the harness's presortedness parameter: $S=2$ is essentially random input, larger $S$ means longer presorted runs.

\begin{table}[h]
\centering
\begin{tabular}{lccc}
\toprule
input & VM-Powersort cmp / mov & SotSort 8-r6 cmp / mov \\
\midrule
\texttt{int}, $S=2$, $n\approx10^6$     & 27.76 / 31.65 & 18.96 / 22.10 \\
\texttt{fastarr}, $S=2$, $n\approx10^6$ & 28.09 / 32.01 & 19.10 / 17.39 \\
\texttt{int}, $S=2$, $n\approx10^7$     & 26.07 / 29.90 & 21.87 / 23.50 \\
\bottomrule
\end{tabular}
\caption{Comparisons and moves per element, same memory.}
\end{table}

Before the paged pass, SotSort's moves were $37.74$ (\texttt{int}, $10^6$) and $48.18$
($10^7$). Across the 12 cells of the harness SotSort is never above VM-Powersort in
comparisons by more than $0.02$ per element. With $S\ge100$ the input takes SotSort's
run-merging path, so the pass is not involved; there, with an $n/2$ buffer,
copy-Powersort still makes fewer moves for $S\ge1000$.

\paragraph{Ablation over $k$} (prototype, 8-byte counted type, VM-Powersort's budget,
random input; moves / comparisons per element).
\begin{center}
\begin{tabular}{lccccc}
\toprule
$n$ & binary (old) & $k=8$ & $k=16$ & $k=32$ & $k=64$\\
\midrule
$10^6$ & 38.46 / 19.61 & 26.76 / 19.69 & 24.19 / 19.64 & 23.81 / 19.59 & 22.95 / 19.64\\
$10^7$ & 48.32 / 22.94 & 31.74 / 23.15 & 29.04 / 23.11 & 26.36 / 23.07 & 26.30 / 23.05\\
\bottomrule
\end{tabular}
\end{center}

\paragraph{Adversary.}
A lazy McIlroy adversary~\cite{mcilroy} run against the prototype, with a fixed seed and
VM-Powersort's budget, at $n=10^6$: with only the depth cap it forced $95.4$
comparisons and $123.3$ moves per element. With the bad-pass rule and the base
case's budget of $4$: $26.13$ comparisons and $34.06$ moves. The constructed input
depends on the seed; with another seed the same input costs $17.37$ / $17.74$.
SotSort accepts a per-call seed. VM-Powersort, as a mergesort, has no such weakness.

\paragraph{Verification.}
The pass was checked exhaustively against \texttt{std::stable\_sort} for all inputs
with keys in $\{0..k-1\}$, $k=2..4$, $P=1..3$, $m\le9$ ($1\,140\,207$ cases), and
the sort on 81 adversarial families and on 112-byte types and \texttt{std::string},
under ASan, UBSan and TSan.

\paragraph{Time.}
Indicative only, ratio to VM-Powersort at $S=2$ on two machines (cloud / desktop):
\texttt{int} 0.72 / 0.78, \texttt{ptr} 0.82 / 0.86, \texttt{fastarr} 0.73 / 0.73,
\texttt{arr} 0.55 / 0.57. Repeated runs vary by about $\pm10$--$15\%$.

\section{Prior work}

The search covered general search engines and the related-work sections of the papers
below; it did not cover patents or all papers citing IPS$^4$o. We may have missed work.

\begin{itemize}
\item \textbf{Polychroniou and Ross}~\cite{pr14}, \S3.2.3, ``out-of-cache, in-place,
list of blocks'': each partition is a linked list of single-destination blocks
written into input space already read. This is the oldest instance of the mechanism we
found. They do not discuss the stability of that variant, state that in-place
partitioning is not stable, and none of their sorts is in-place, stable and
block-based at once.
\item \textbf{Radsort} (Clausecker and Schintke)~\cite{radsort}: a stable LSD radix sort that
reuses consumed blocks and tracks the block permutation in arrays, explicitly
comparing its design to Polychroniou and Ross; an MSD version is left as future
work. This is the closest prior work; our pass uses the same core mechanism.
\item \textbf{IPS$^4$o}~\cite{ips4o}: flushes full buffer blocks into the input but,
to stay in-place, does not record which bucket each flushed block belongs to;
it is unstable. Recording the group (with $O(n/P)$ words) is precisely what buys stability.
\item \textbf{Franceschini, Muthukrishnan and P\u{a}tra\c{s}cu}~\cite{fmp07}: merging $k$ sorted
lists with $\alpha n$ free words by writing output into free blocks and reclaiming blocks freed
in the inputs: the same reuse idea, for merging.
\item \textbf{Katajainen and Pasanen}~\cite{kp92,kp94}: stable partitioning and stable
multiset sorting in minimum space, quicksort with median pivots; theoretical.
\item \textbf{Logsort}~\cite{logsort}: a practical stable in-place 2-way quicksort.
\item External-memory distribution sort writes each bucket in blocks~\cite{vitter};
it does not reuse the input array.
\end{itemize}

What we claim as new is narrow: a comparison-based MSD samplesort whose distribution
pass is stable, reuses read pages, and leaves groups contiguous after each pass, with a
bounded buffer. We have not found it, and S.~Wild did not know of one offhand.

\section*{Acknowledgements}
Sebastian Wild advised focusing on robust algorithmic ideas and operation counts rather
than running time. Errors are mine.

\begin{thebibliography}{9}
\bibitem{ips4o} M.~Axtmann, S.~Witt, D.~Ferizovic, P.~Sanders. Engineering in-place
(shared-memory) sorting algorithms. \emph{ACM Trans.\ Parallel Comput.} 9(1), 2022.
Earlier version: In-place parallel super scalar samplesort (IPS$^4$o), ESA 2017.
\bibitem{radsort} R.~Clausecker, F.~Schintke. Parallel $O(\sqrt n)$ overhead LSD radix sort.
arXiv:2607.05302, 2026.
\bibitem{pr14} O.~Polychroniou, K.~A.~Ross. A comprehensive study of main-memory partitioning
and its application to large-scale comparison- and radix-sort. SIGMOD 2014, 755--766.
\bibitem{fmp07} G.~Franceschini, S.~Muthukrishnan, M.~P\u{a}tra\c{s}cu. Radix sorting with no
extra space. ESA 2007, LNCS 4698. arXiv:0706.4107.
\bibitem{kp92} J.~Katajainen, T.~Pasanen. Stable minimum space partitioning in linear time.
\emph{BIT} 32(4):580--585, 1992.
\bibitem{kp94} J.~Katajainen, T.~Pasanen. Sorting multisets stably in minimum space.
\emph{Acta Informatica} 31(4):301--313, 1994.
\bibitem{logsort} aphitorite. Logsort. \url{https://github.com/aphitorite/Logsort}.
\bibitem{vitter} J.~S.~Vitter. External memory algorithms and data structures: dealing with
massive data. \emph{ACM Comput.\ Surv.} 33(2), 2001.
\bibitem{vmps} F.~Moltmann, T.-V.~Nakajima, S.~Wild. Virtual-Memory Powersort. ESA 2026,
LIPIcs 388, 14:1--14:17. Harness use as in L.~Sotomayor, SotSort, arXiv:2611.24786.
\bibitem{mcilroy} M.~D.~McIlroy. A killer adversary for quicksort. \emph{Softw.\ Pract.\ Exper.}
29(4):341--344, 1999.
\end{thebibliography}
\end{document}
